% Preámbulo de la colección de problemas.
% Edítalo a tu gusto; debe terminar justo después de \begin{document}
% (y, si quieres, \maketitle / \tableofcontents). El generador añade el cuerpo
% con los \input de cada problema y cierra con \end{document}.
\documentclass[11pt]{article}
\usepackage[utf8]{inputenc}
\usepackage[T1]{fontenc}
\usepackage[spanish,es-noshorthands]{babel}
\usepackage{amsmath}
\usepackage{amssymb}
\usepackage{graphicx}
\usepackage{comment}
\usepackage{booktabs}
\usepackage{multirow}
\usepackage{enumitem}
\usepackage{float}
\usepackage{tikz}
\usetikzlibrary{arrows, arrows.meta, positioning, calc, automata, graphs, shapes.geometric}
% --- Estilos TikZ personalizados ---------------------------------------
% Algunos problemas de grafos usan estilos propios (p. ej. arn_n). Pega aquí
% las definiciones de tu preámbulo original para que se rendericen bien:
% \tikzstyle{arn_n}=[circle, fill=black!90, text=white, minimum size=22pt]
% -----------------------------------------------------------------------
\usepackage[margin=2.5cm]{geometry}

\title{Colección de problemas de Programación Lineal}
\author{}
\date{}

\begin{document}
\maketitle
\tableofcontents
\clearpage
